\chapter{CONCLUSIONS AND FUTURE SCOPE}

\paragraph{}
\textbf{Conclusion.} The project turned a working but fragile one-way pipeline into a Lambda architecture that treats data loss as the primary design problem. MongoDB is now sharded (hashed key, two shards) and replicated (three members per shard), every event is also stored in an immutable Parquet lake, snapshots are verified and rotated, and streaming checkpoints are archived with them. Three destructive tests were demonstrated and measured: a crashed shard primary, a dropped database, and a dropped database with all backups deleted; none lost a document, and recovery took between 1 and 10 seconds. A Spark MLlib payload model raised recall on the full dataset from 67.67\% to 99.85\% with the label never used at inference, a four-partition Kafka topic improved ingestion by 34\%, and a Spark SQL batch layer and an extended dashboard complete the architecture. 106 automated tests pass. The work applies sharding, replication, columnar storage, batch versus stream processing and window aggregation from the Big Data Analytics course to a concrete problem.

\paragraph{}
\textbf{Limitations.} Everything runs on one laptop, so the cluster shows correct distribution and failover but no throughput gain and no protection against loss of the machine; the 20,000 records/s target was not met; the false-positive rate rose to 5.15\%; the ML model was trained on a small sample from a highly separable dataset; and recovery from the lake cannot restore per-rule alert reasons or original batch timestamps.

\paragraph{}
\textbf{Future scope.} (1) Place shards, replicas and the backup target on separate machines or an object store for a true 3-2-1 backup policy and to realise throughput scaling. (2) Add point-in-time recovery from the replication log between snapshots. (3) Replicate Kafka itself (three brokers, replication factor three) so the stream survives a broker loss. (4) Tune the hybrid threshold or use a calibrated stacking model to reduce false positives, and train on a larger sample. (5) Ingest live packet captures instead of a replay, and add per-alert severity. (6) Evaluate GPU-accelerated gradient boosting for the payload model.
